\documentclass[11pt]{article}
\usepackage[a4paper,margin=25mm]{geometry}
\usepackage[T1]{fontenc}
\usepackage[utf8]{inputenc}
\usepackage{amsmath,amssymb,amsthm}
\usepackage{hyperref}
\hypersetup{hidelinks}
\usepackage{listings}
\lstset{basicstyle=\ttfamily\small,breaklines=true,columns=fullflexible}
\newtheorem{theorem}{Theorem}
\newtheorem{lemma}[theorem]{Lemma}
\newcommand{\sfcount}{S}
\newcommand{\mm}{m}
\title{Conjecture 00000001337: the iterated M\"obius mean\\converges to $6/\pi^2$ with an explicit error}
\author{AlyciaBHZ on behalf of the Omega Institute}
\date{5 October 2026}
\begin{document}
\maketitle
\begin{abstract}
With the convention specified in the conjecture, the composition
$\mu(\mu(n))$ is exactly the indicator of squarefree positive integers.
We prove, for every real $x\geq0$,
\[
\left|\sum_{1\leq n\leq x}\mu(\mu(n)) - \frac{6}{\pi^2}x\right|
\leq 3\sqrt{x}.
\]
Consequently its mean converges to $6/\pi^2>0$ and cannot converge to zero.
The accompanying self-contained Lean~4 project proves the literal
composition identity, the error bound, the real-endpoint limit, and the
failure of convergence to zero. All audited theorems use only the standard
axioms \texttt{propext}, \texttt{Classical.choice}, and \texttt{Quot.sound}.
\end{abstract}

\section{Statement and interpretation}
The statement of problem \textbf{00000001337}, verbatim (with
mathematical symbols typeset), is:
\begin{quote}
``Definition: Applying $\mu$ to the values of $\mu(n)$, i.e.
$\mu(\mu(n))$ (with $\mu(n)\in\{-1,0,1\}$, setting $\mu(\pm1)=1$
and $\mu(0)=0$). Conjecture: The mean
$(1/x)\Sigma_{n\leq x}\,\mu(\mu(n))$ does not converge to 0 but diverges
to a $6/\pi^2$-type limit --- the value statistics of $\mu$ interact with
its zero density to produce a main term.''
\end{quote}

The usual M\"obius function on positive integers is
\[
\mu(n)=\begin{cases}
0,&\text{if a prime square divides }n,\\
(-1)^r,&\text{if }n\text{ is a product of }r\text{ distinct primes}.
\end{cases}
\]
In particular $\mu(1)=1$. We interpret the summation as over positive
integers, as is conventional for this arithmetic mean. For the outer
application, define on \emph{all} integers
\[
\mu_{\mathrm{out}}(z)=
\begin{cases}1,&z=1\text{ or }z=-1,\\0,&\text{otherwise},\end{cases}
\qquad \mm(n)=\mu_{\mathrm{out}}(\mu(n)).
\]
This explicitly encodes all three stated outer values. Assigning zero to
other integer inputs is immaterial because the inner function takes only
$-1,0,1$. We do not interpret the outer function as the usual M\"obius
function applied to a signed integer without the stated extension.

The phrase ``diverges to a limit'' is internally inconsistent if
``diverges'' has its standard meaning of failure of convergence. The
faithful precise claim is that this mean \emph{converges} to $6/\pi^2$,
and thus does not tend to zero. We prove the exact constant, rather than
only a vaguely similar ``type'' of limit. If divergence is insisted upon
literally, that additional assertion is false: our convergence theorem
settles the actual behavior. The nonzero-limit part of the conjecture is
true.

\section{The composition and a finite counting identity}
A positive integer is squarefree if it has no prime-square divisor.
The values of the ordinary M\"obius function immediately give
\begin{equation}\label{eq:indicator}
\mm(n)=\begin{cases}1,&n\text{ squarefree},\\0,&n\text{ not squarefree}.
\end{cases}
\end{equation}
Write
\[
\sfcount(x)=\sum_{1\leq n\leq x}\mm(n),\qquad
M=\lfloor\sqrt{x}\rfloor\quad(x\geq0).
\]
Both sums are empty when their upper endpoint is below $1$.

\begin{lemma}\label{lem:divisor}
For each positive integer $n$,
\[
\mathbf 1_{\mathrm{squarefree}}(n)=\sum_{d^2\mid n}\mu(d).
\]
Consequently
\begin{equation}\label{eq:count}
\sfcount(x)=\sum_{1\leq d\leq M}\mu(d)
             \left\lfloor\frac{x}{d^2}\right\rfloor
\qquad(x\geq0).
\end{equation}
\end{lemma}
\begin{proof}
Factor $n=\prod_p p^{a_p}$. The only nonzero contributions to the divisor
sum come from squarefree $d$, made from primes with $a_p\geq2$.
The sum is therefore
$\prod_{p:a_p\geq2}(1+\mu(p))=\prod_{p:a_p\geq2}(1-1)$.
It equals $1$ when that product is empty, and $0$ otherwise, proving the
indicator identity. Equivalently, if $n=t^2b$ with $b$ squarefree and
$t=\prod_p p^{\lfloor a_p/2\rfloor}$, then $d^2\mid n$ if and only if
$d\mid t$; the divisor sum is $\sum_{d\mid t}\mu(d)=\mathbf1_{t=1}$.

Summing the identity over $1\leq n\leq\lfloor x\rfloor$ and interchanging
finite sums, only $d\leq\sqrt{x}$ occur. For fixed $d$, the number of
positive multiples of $d^2$ not exceeding $x$ is
$\lfloor x/d^2\rfloor$. This proves \eqref{eq:count} for real as well as
integer $x$, including $x=0$.
\end{proof}

\section{The coefficient and the tail}
\begin{lemma}\label{lem:coefficient}
The absolutely convergent series satisfies
\[
c:=\sum_{d=1}^{\infty}\frac{\mu(d)}{d^2}=\frac6{\pi^2}.
\]
\end{lemma}
\begin{proof}
Since $|\mu(d)|\leq1$ and $\sum_{d\geq1}d^{-2}=\pi^2/6$,
the M\"obius series converges absolutely. Its product with the latter
positive series may be rearranged by absolute convergence:
\[
\left(\sum_{d\geq1}\frac{\mu(d)}{d^2}\right)
\left(\sum_{k\geq1}\frac1{k^2}\right)
=\sum_{n\geq1}\frac1{n^2}\sum_{d\mid n}\mu(d)=1.
\]
To justify the last identity directly, the inner divisor sum equals
$\prod_{p\mid n}(1-1)$ for $n>1$, and equals $1$ for $n=1$.
The double absolute sum is at most
$(\sum_{d\geq1}d^{-2})^2<\infty$, which justifies all regrouping.
The Basel sum evaluates the second factor as $\pi^2/6$, proving the claim.
\end{proof}

\begin{lemma}\label{lem:tail}
For every real $N\geq1$,
\[
\sum_{j=0}^{\infty}\frac1{(j+N)^2}\leq\frac2N.
\]
In particular
\[
\left|\sum_{d>M}\frac{\mu(d)}{d^2}\right|\leq\frac2{M+1}
\quad\text{for every integer }M\geq0.
\]
\end{lemma}
\begin{proof}
The function $t\mapsto(t+N)^{-2}$ is nonnegative and decreasing for
$t\geq0$. Each term with $j\geq1$ is at most the integral over
$[j-1,j]$, so
\[
\sum_{j\geq0}\frac1{(j+N)^2}
\leq\frac1{N^2}+\int_0^\infty\frac{dt}{(t+N)^2}
=\frac1{N^2}+\frac1N\leq\frac2N.
\]
Absolute convergence and $|\mu(d)|\leq1$ give the second assertion by
setting $N=M+1$.
\end{proof}

\section{The explicit error and convergence}
\begin{theorem}\label{thm:error}
For every real $x\geq0$,
\begin{equation}\label{eq:error}
\left|\sfcount(x)-\frac6{\pi^2}x\right|\leq3\sqrt{x}.
\end{equation}
For every real $x>0$,
\begin{equation}\label{eq:meanerror}
\left|\frac{\sfcount(x)}x-\frac6{\pi^2}\right|
\leq\frac3{\sqrt{x}}.
\end{equation}
\end{theorem}
\begin{proof}
Set $M=\lfloor\sqrt{x}\rfloor$ and use Lemmas~\ref{lem:divisor}
and~\ref{lem:coefficient}. Subtracting the infinite series coefficient,
\[
\sfcount(x)-cx
=\sum_{d=1}^{M}\mu(d)
  \left(\left\lfloor\frac{x}{d^2}\right\rfloor-\frac{x}{d^2}\right)
 -x\sum_{d>M}\frac{\mu(d)}{d^2}.
\]
The absolute value of the finite sum is at most $M$, because each
floor error has absolute value at most $1$ and $|\mu(d)|\leq1$.
Lemma~\ref{lem:tail} bounds the remaining part by $2x/(M+1)$.
Since
\[
M\leq\sqrt{x}\leq M+1,\qquad x=(\sqrt{x})^2,
\]
we have $x/(M+1)\leq\sqrt{x}$. Hence
\[
|\sfcount(x)-cx|\leq M+\frac{2x}{M+1}\leq3\sqrt{x}.
\]
This argument includes $M=0$ and $x=0$. Dividing by $x>0$ gives
\eqref{eq:meanerror}, since $\sqrt{x}/x=1/\sqrt{x}$.
\end{proof}

\begin{theorem}\label{thm:limit}
As the real variable $x\to+\infty$,
\[
\frac1x\sum_{1\leq n\leq x}\mu_{\mathrm{out}}(\mu(n))
\longrightarrow\frac6{\pi^2}>0.
\]
In particular this mean does not converge to zero.
\end{theorem}
\begin{proof}
The right side of \eqref{eq:meanerror} tends to zero, so the squeeze
principle proves the stated limit. Since $\pi>0$, we have
$6/\pi^2>0$. Uniqueness of limits in the real numbers excludes a second
limit equal to zero.
\end{proof}
Thus the zero density of the inner M\"obius function is responsible for
the missing mass; the outer convention removes its signs. The main term
is exactly the squarefree density, not a signed cancellation effect.

\section{Lean statements and semantic correspondence}
The project pins Lean~\texttt{v4.33.0} and Mathlib~\texttt{v4.33.0}.\\
Mathlib commit: \texttt{db584cd6d46c92f209a44c0f1c829460d327499d}.

The single source file is \texttt{lean4/MobiusMean.lean}. The declarations
are in the namespace \texttt{MobiusMean}. Written with equivalent ASCII Lean syntax,
its definitions are:
\begin{lstlisting}
def mobiusOuter (z : Int) : Int :=
  if Or (z = 1) (z = -1) then 1 else 0
noncomputable def mm (n : Nat) : Int :=
  mobiusOuter (ArithmeticFunction.moebius n)
noncomputable def count (x : Real) : Real :=
  Finset.sum (Finset.Icc 1 (Nat.floor x)) (fun n => (mm n : Real))
noncomputable def mean (x : Real) : Real := (1 / x) * count x
\end{lstlisting}
Here \texttt{Finset.sum} writes out the sum notation used in the delivered
source, and \texttt{Nat}, \texttt{Int}, \texttt{Real} are the types
usually written with blackboard-bold symbols.
For $x\geq0$, \texttt{Nat.floor x} is $\lfloor x\rfloor$; its convention
at negative $x$ is harmless for a limit at positive infinity.
The inclusive interval starts at $1$, so the summation is precisely over
positive integers at most $x$.

The central theorem statement in equivalent ASCII Lean syntax is:
\begin{lstlisting}
theorem mean_tendsto :
  Tendsto (fun x : Real => (1 / x) *
    Finset.sum (Finset.Icc 1 (Nat.floor x)) (fun n => (mm n : Real)))
    atTop (nhds (6 / Real.pi ^ 2))
\end{lstlisting}
The same function appears in \texttt{not\_tendsto\_zero}, whose
conclusion is the negation of \texttt{Tendsto ... atTop (nhds 0)}.
The remaining main statements are:
\begin{lstlisting}
theorem mm_eq_indicator (n : Nat) :
  mm n = if Squarefree n then 1 else 0

theorem count_error (x : Real) (hx : 0 <= x) :
  abs (Finset.sum (Finset.Icc 1 (Nat.floor x))
    (fun n => (mm n : Real))
    - (6 / Real.pi ^ 2) * x) <= 3 * Real.sqrt x

theorem mean_error (x : Real) (hx : 0 < x) :
  abs (mean x - 6 / Real.pi ^ 2) <= 3 / Real.sqrt x

theorem limit_ne_zero : Ne (6 / Real.pi ^ 2 : Real) 0
\end{lstlisting}
These statements have no conditional assumption of a density, an
unproved asymptotic, or a numerical approximation to $\pi$.
The only hypotheses of the error estimates restrict the real endpoint
as stated in Theorem~\ref{thm:error}; the limit is unconditional.

The supporting vendored theorem actually proves the more general
squarefree count coprime to a squarefree modulus $Q$:
\[
|S_Q(x)-\rho(Q)x|\leq(2^{\omega(Q)}+2)\sqrt{x},\qquad
\rho(Q)=\left(\sum_{n\geq0}\frac1{(n+1)^2}\right)^{-1}
\prod_{p\mid Q}\frac p{p+1}.
\]
In the source, \texttt{S} is a filtered cardinality and
\texttt{rho} is exactly this expression. Its proof is included in full;
it is not assumed. At $Q=1$ all integers are coprime to $Q$, the prime
factor set is empty, and $2^{\omega(1)}+2=3$.
\texttt{count\_eq\_S} identifies this cardinality with the sum of the
literal composition. \texttt{rho\_one} uses Mathlib's
\texttt{hasSum\_zeta\_two} to prove $\rho(1)=6/\pi^2$.
The vendored coefficient proof uses the M\"obius inverse identity for
the principal Dirichlet-character $L$-series at $s=2$; its specialization
is exactly the absolutely convergent product in
Lemma~\ref{lem:coefficient}. The floor errors and the tail estimate are
proved inside this source. The general coprime result is auxiliary;
the mathematical proof above gives all arguments needed for the submitted
$Q=1$ conclusion.

\section{Verification and axiom audit}
The actual compiler transcript is in \texttt{verification.txt}.
The recorded Lean compilation succeeded with the pinned Lean and Mathlib versions.
The source ends with \texttt{\#print axioms} for each of
\texttt{mm\_eq\_indicator}, \texttt{count\_error},
\texttt{mean\_error}, \texttt{mean\_tendsto},
\texttt{limit\_ne\_zero}, and \texttt{not\_tendsto\_zero}.
Every one reports exactly
\begin{lstlisting}
[propext, Classical.choice, Quot.sound]
\end{lstlisting}
No new axioms, admitted goals, or native decision procedure are used.
The proof is checked by Lean's kernel. Noncomputable definitions use
standard classical choice and do not introduce any unproved assertion.

The independent standard-library script \texttt{verify.py} checks the
outer convention, compares M\"obius values with trial factorization,
verifies the exact finite floor-sum identity for $0\leq N\leq500$, and
checks both error estimates at every integer endpoint up to $200000$
and seven additional nonintegral endpoints. For example,
$S(200000)=121581$ and the mean is $0.607905$, while
$6/\pi^2\approx0.607927101854$. These floating-point checks are sanity
checks only; the Lean statements prove all endpoints and the infinite
limit without numerical approximation.

\section{Other readings, provenance, and reproduction}
For integer cutoffs, the same limit follows by restriction of the proved
real-endpoint limit to $x=N\in\mathbb N$ tending to infinity; the error
bound already applies to every such $N$. Replacing the normalizing $x$
by $\lfloor x\rfloor$ for $x\geq1$ also gives the same limit: the numerator
is $S(\lfloor x\rfloor)$ and $\lfloor x\rfloor/x\to1$.
Using $n<x$ instead of $n\leq x$ changes the numerator by at most one
term, hence the mean by at most $1/x$. Including $n=0$, with
$\mu(0)=0$, adds zero. These variants do not change the conclusion;
the exact statements formalized in Lean use inclusive cutoffs and the
normalization $1/x$.

The proof of the supporting coprime squarefree count, the tail lemma, and
the definitions \texttt{C}, \texttt{e} are adapted from the Omega Institute's
\emph{trureturing} project, licensed under Apache~2.0, at revision
\texttt{1be30b25c4fe6cd46397090ca09b5696eb62ea78}.
The credited source paths are
\begin{itemize}
\item \texttt{D5/S3/Weil/Mertens/CoprimeSquarefreeDensity.lean};
\item \texttt{D5/S3/Weil/Mertens/Third.lean} (only
\texttt{sum\_one\_div\_sq\_le});
\item \texttt{D5/S3/Weil/Mertens/CoprimeMobiusCertificateError.lean}
(only \texttt{C} and \texttt{e}).
\end{itemize}
All code needed from these files is inlined and renamed into the
problem-specific namespace. There are no trureturing imports or external
proof certificates. The upstream license accompanies the submission as
\texttt{LICENSE.trureturing}.

In a fresh environment with the pinned toolchain, reproduce with:
\begin{lstlisting}
cd lean4
lake exe cache get && lake build
cd ..
python3 verify.py
pdflatex -interaction=nonstopmode -halt-on-error report.tex
pdflatex -interaction=nonstopmode -halt-on-error report.tex
\end{lstlisting}

\medskip
\noindent Submitted by AlyciaBHZ on behalf of the Omega Institute
(trureturing project,
\url{https://github.com/the-omega-institute/trureturing}).
\end{document}
